\documentclass[11pt,a4paper]{report}
\usepackage[utf8]{inputenc}
\usepackage[T1]{fontenc}
\usepackage[margin=2.5cm]{geometry}
\usepackage{hyperref}
\renewcommand{\contentsname}{Índex}
\renewcommand{\chaptername}{Capítol}
\title{Projecte Intermodular SIMIX2\\Xarxa de l'Institut}
\author{Joan Ventura \and Lucas Marés \and Aleix Mugica}
\date{Curs 2026--2027 · SMX2}
\begin{document}
\maketitle
\tableofcontents
\chapter{Presentació del projecte}
\section{Nom i lema}
El projecte s'anomena \textbf{Xarxa de l'Institut}. El lema és: «Un institut amb una xarxa que funciona de veritat... o almenys simulada».
\section{Problema que volem resoldre}
Un institut necessita connectar els ordinadors de secretaria, professorat i alumnat, compartir documents i oferir informació a les famílies. Cada grup ha de poder fer la seva feina sense accedir als fitxers que no li corresponen. També cal poder recuperar la informació si hi ha una avaria.

El projecte proposa dissenyar i provar aquesta xarxa en un entorn de pràctiques. Les proves documentades a la sessió 3 s'han fet amb Packet Tracer; no demostren encara el funcionament dels mateixos serveis en un servidor real.
\section{Persones usuàries}
\begin{itemize}
\item Secretaria: treballar amb els equips de gestió del centre.
\item Professorat: compartir materials i comunicar-se amb el personal.
\item Alumnat: consultar i descarregar els materials autoritzats.
\item Famílies: consultar la informació pública del centre al web.
\item Administració de la xarxa: mantenir els equips, els accessos i les còpies.
\end{itemize}
\section{Objectius}
\begin{enumerate}
\item Separar les zones de secretaria, professorat, aules i servidors.
\item Assignar automàticament les adreces als equips i permetre trobar els serveis pel seu nom.
\item Preparar un espai de fitxers compartits amb permisos diferents per a professorat i alumnat.
\item Preparar correu intern i manteniment a distància amb accés restringit.
\item Crear un web informatiu amb VS Code, sense WordPress.
\item Preparar la connexió sense fils, controlar els accessos i provar la recuperació de còpies.
\item Guardar les proves i explicar els resultats perquè un altre company els pugui revisar.
\end{enumerate}
\section{Abast inicial}
El disseny de la sessió 2 preveu 17 ordinadors: 3 a secretaria, 4 per al professorat i 10 a les aules. Les quatre zones es distingeixen amb els números 10, 20, 30 i 40. Els servidors tenen una adreça fixa i els equips de les tres primeres zones reben una adreça automàtica.

La primera versió del web oferirà informació pública. Qualsevol funció amb dades personals, com ara les notes de l'alumnat, requerirà una definició i una revisió separades; no s'inclou com a funció ja resolta.
\section{Equip i repartiment}
\begin{description}
\item[Joan Ventura] Xarxa i equips; suport a la documentació.
\item[Lucas Marés] Serveis de xarxa; suport a la seguretat.
\item[Aleix Mugica] Web; suport a la documentació.
\end{description}
\section{Situació documentada}
Els Word S0 i S1 recullen la idea, les necessitats, les prioritats i la planificació. El S2 recull el disseny de la xarxa. El S3 aporta captures de l'assignació d'adreces i la resolució de noms a Packet Tracer, i una reunió datada el 29 de setembre de 2026.

Aquests documents descriuen sessions de setembre. La pauta de la primera entrega fixa el Sprint 1 del 6 al 20 d'octubre de 2026. El seguiment setmanal i les reunions finals d'aquest període s'han d'afegir amb les dates i els resultats reals. No es donen per fets en aquesta memòria.
\section{Fonts i seguiment}
Aquesta versió inicial s'ha redactat a partir dels quatre Word aportats per l'equip i de la presentació del repositori. Els originals es conserven a la carpeta docs.

Repositori: \url{https://github.com/joanvent/projecteintermodular}

Tauler: \url{https://github.com/users/joanvent/projects/2}
\end{document}
